\section{Conclusions and Extensions}\label{sec:conclusion}
This paper developed a prescribed-time controller based on smooth periodic delayed feedback for joint-space tracking of a pre-capture free-floating space manipulator. The four-stage structure combines disturbance observation, dynamic compensation, error linearisation, and periodic delayed feedback. Under continuous-time implementation, an exact model, no saturation, and exact compensation, tracking errors vanish after $T_{\mathrm{pre}}=T_o+2h$; with residual disturbances and approximate nullification, explicit practical bounds were established. The design was then evaluated on the complete spatial model of a seven-degree-of-freedom manipulator parameterised with the data of~\cite{dou2025fractional}, using a composite-rigid-body/recursive-Newton--Euler implementation, the momentum-consistent base reaction, and the tabulated $\pm20~\mathrm{N\,m}$ actuator limits.

The numerical study supports three conclusions. First, the one-period nullification is observable in a fully coupled three-dimensional plant: after the prescribed time the residual error of the delayed-feedback law is $4.3\times10^{-5}~\mathrm{rad}$, more than one order of magnitude below the proportional--derivative baseline, and the measured settling time is about half of the baseline value. Second, the observer, not the delayed feedback, is responsible for disturbance rejection: it reduces the tracking error from $2.6\times10^{-1}~\mathrm{rad}$ to $1.5\times10^{-3}~\mathrm{rad}$ within the actuator envelope, which confirms the role of the prescribed-time observer as the interface between the uncertain plant and the delayed-feedback error channel. Third, the equivalent inertia of the arm spans three decades, and this ratio, rather than the control law, determines where the design margins are: it fixes the admissible initial error for a given torque limit, it sets the numerical conditioning of the Gramian, and it makes a single delay ill-matched to the fastest joint modes.

These observations suggest several concrete extensions.

\begin{enumerate}
\item \emph{Structured and multi-rate delayed feedback.} The present design uses one delay $h$ and one activation function for all joints. A per-channel delay, or a delay matrix built from the dominant time constant of each joint group, would let the slow shoulder joints and the fast wrist joints be nullified in their own time scales. The algebra of Theorem~\ref{thm:pdf_nullification} extends to a block-diagonal gain with different delays if the period-boundary operator is evaluated blockwise, and the condition number of the resulting Gramian would become a design objective rather than a numerical accident.
\item \emph{Saturation-aware nullification.} The $3^\circ$ case shows that the delayed feedback loses its one-period property when the actuator limit is active for a non-negligible fraction of the design window. Two complementary routes seem promising: an event-triggered re-nullification that restarts the window when the error re-enters the actuator envelope, and a feasibility condition that maps the triple (initial error, torque limit, delay) into an admissible set. The fixed-time anti-saturation compensator of~\cite{dou2025fractional} and the fragility-avoidance prescribed-performance design of~\cite{zhou2025saturated} provide ready-made components for the second route.
\item \emph{Task-space and base-attitude requirements.} The law is formulated in joint space; a complete on-orbit servicing controller must also keep the end effector on a Cartesian path while limiting the base attitude drift. A natural next step is to close the loop through the generalized Jacobian with a damped pseudo-inverse, to use the reaction null space for attitude regulation~\cite{zong2020reactionless,xu2016}, and to check the reference against the dynamic-singularity loci during planning~\cite{nanos2015avoiding,nanos2023fullpose}. In that setting the prescribed-time guarantee would apply to the joint reference generated by the task-space layer, and the task-space residual would enter the practical bound explicitly.
\item \emph{Nonzero momentum and free-flying operation.} Assumption~1 restricts the analysis to zero initial momentum. When reaction wheels or thrusters are present, the base velocity is no longer an algebraic function of the joint rates, the reduced model acquires an additional input channel, and the delay design must nullify a larger state. The augmented-momentum formulations of~\cite{nanos2017momentum,nanos2023fullpose} provide the modelling framework, and the momentum distribution between base and arm becomes an additional degree of freedom that the delayed feedback could exploit.
\item \emph{Sampled-data and quantised implementation.} The residual error at the prescribed time in the reported runs is set by the explicit Euler step and by the sampled delay buffer, not by the continuous-time design. An exact discretisation of the delay system, together with a discrete-time version of the one-period map, would replace the empirical scaling with a certified bound; input quantisation and measurement noise can be treated in the same framework~\cite{hu2025preassigned}.
\item \emph{Richer plant models.} Flexibility of the appendages, joint friction, actuator dynamics, and contact with the target are all absent from the present model. They matter for a pre-capture manoeuvre because they introduce unmodelled high-frequency dynamics and, at contact, a discontinuous change of the plant. Extending the practical bound to a plant with a flexible mode, or switching the delayed feedback off during contact and restarting it afterwards, are well-defined next problems.
\item \emph{Experimental validation and benchmarking.} The parameter set, the actuator limits, the settling-time definition, the peak torque, and the saturation duty cycle are all reported here, which makes the study reproducible and comparable with the sliding-mode results obtained for the same manipulator~\cite{dou2025fractional}. Air-bearing-table experiments, or a flight demonstration in the spirit of~\cite{yoshida2003etsvii}, would be the decisive test, since the delay buffer and the actuator limit are the two elements most sensitive to hardware.
\end{enumerate}

Finally, the results suggest a broader theoretical question: the delayed-feedback construction converts a convergence-time specification into a numerical linear-algebra problem (the annihilation of a period map), and the practical bound converts the modelling error into the conditioning of that problem. Making this trade-off explicit --- for example, by minimising the Gramian condition number subject to a torque budget --- appears to be a fruitful route towards prescribed-time control that is simultaneously rigorous and implementable on torque-limited space manipulators.
